\chapter{Alcance y delimitaciones}

El proyecto abordó la reactivación progresiva del robot cuadrúpedo Nova Spot
Micro mediante modelado, simulación, control de locomoción, preparación
electrónica y entrenamiento y evaluación de una capa correctiva de aprendizaje por refuerzo.
El sistema se limitó a postura y locomoción de baja velocidad tipo gateo, paso
y trote, evaluadas en simulación y en el robot físico suspendido, bajo
condiciones controladas.

\section{Componentes desarrollados}

\begin{itemize}
\item Se formularon e implementaron modelos cinemáticos y dinámicos nominales,
jacobianos y descripciones URDF/SDF y MJCF.
\item Se integró una arquitectura ROS~2 con control articular, métricas,
supervisor, fase de marcha, contacto, IMU y margen de estabilidad nominal.
\item Se implementaron y evaluaron en simulación postura, marcha paso y gateo.
El galope se conservó únicamente como experimento opcional, deshabilitado por
defecto y excluido de los criterios principales.
\item Se preparó la Raspberry Pi como controlador embebido, se verificó
comunicación DDS e I$^2$C y se realizaron pruebas progresivas del PCA9685.
\item Se instrumentaron las articulaciones con potenciómetros y encoders
magnéticos, se identificó el actuador real y se ejecutaron marchas nominales y
políticas PPO en el robot suspendido, con registro sincronizado de sensores y
hasta cinco cámaras.
\item Se entrenaron y seleccionaron políticas PPO residuales con el actuador
identificado, mediante un ciclo que evalúa las candidatas en el robot y
actualiza el modelo con sus datos.
\item Se definieron el protocolo, las métricas, las cantidades experimentales y
las semillas para la comparación final.
\end{itemize}

\section{Delimitaciones}

No se abordaron navegación autónoma, conectividad en la nube, percepción del
entorno ni locomoción rápida sobre terreno irregular. La política aprendida no
comandará directamente PWM ni sustituirá límites articulares, saturaciones o el
supervisor independiente.

Se identificaron la geometría de las patas y el actuador (retardo, zona
muerta, histéresis y saturación), pero no las masas, inercias ni la fricción, por
lo que el modelo sigue considerándose nominal computable y no un gemelo
digital. El margen de estabilidad calculado en simulación no se presentó como
medición física. El seguimiento articular real se midió con los sensores
añadidos, salvo en dos tibias sin sensor operativo, que se estimaron con la
orden.

La primera estrategia PPO, entrenada con el modelo nominal, no demostró la
mejora requerida y no fue aceptada para transferencia. Las estrategias
posteriores, entrenadas con el actuador identificado, sí se ejecutaron en el
robot suspendido. La seguridad eléctrica y la alimentación (separación entre
electrónica y servos) siguen abiertas para OE1. OE3 requiere demostrar estabilidad y continuidad del control aprendido
en paso y gateo. OE4 exige completar la Prueba de Margen de Estabilidad Estática,
la Prueba de Seguimiento de Posición Articular y la Prueba de Repetibilidad del
Patrón de Marcha; la evidencia parcial en simulación no cierra su evaluación integral.

\section{Productos obtenidos y productos abiertos}

Se obtuvieron código versionado, modelos de simulación, documentación
matemática, controladores nominales, pruebas automatizadas, bolsas de datos,
analizadores, informes, un protocolo experimental, un conjunto de datos físicos
(caminatas con sensores y video sincronizados), políticas PPO seleccionadas en
el robot y su comparación con la marcha nominal. Quedan abiertas la validación
de seguridad eléctrica, la locomoción apoyada con medición de la orientación
del cuerpo y una comparación estadística con más repeticiones. La separación evita presentar productos esperados como si
hubieran sido obtenidos.
